\documentclass{article}
% Packages
\usepackage[utf8]{inputenc}
\usepackage{geometry}
\usepackage[T1]{fontenc}
\usepackage{amsfonts}
\usepackage{amsmath}
\usepackage{graphicx}
\usepackage{float}
\usepackage{hyperref}
\usepackage[sorting=none]{biblatex}
\usepackage{fancyhdr}
\usepackage{multicol}
\usepackage{xcolor}
\usepackage{listings}
\usepackage{tabularx}
\usepackage{rotating}
\renewcommand{\arraystretch}{1.2}

\usepackage{tikz}
\usetikzlibrary {automata, shapes}
\usetikzlibrary{arrows.meta,positioning}
\usepackage{dot2texi}

\addbibresource{ref.bib}
\setlength{\columnsep}{40pt}
\setlength{\voffset}{0.7cm}
\setlength{\headsep}{40pt}
\geometry{legalpaper, portrait, margin=2cm}

\definecolor{codegreen}{rgb}{0,0.6,0}
\definecolor{codegray}{rgb}{0.7,0.7,0.7}
\definecolor{codeorange}{rgb}{0.92,0,0.82}
\definecolor{backcolour}{rgb}{0.91,0.93,0.95} % Light blue-gray for blocks
\definecolor{inlinebackcolour}{rgb}{0.85,0.88,0.91} % Slightly darker for inline

\lstdefinestyle{mystyle}{
    backgroundcolor=\color{backcolour},    
    commentstyle=\color{codegray},
    keywordstyle=\color{blue},
    numberstyle=\tiny\color{codegray},
    stringstyle=\color{codeorange},
    basicstyle=\ttfamily\footnotesize,
    breakatwhitespace=false,          
    breaklines=true,                  
    captionpos=b,                     
    keepspaces=true,                  
    numbers=left,                     
    numbersep=5pt,                    
    showspaces=false,                 
    showstringspaces=false,
    showtabs=false,                   
    tabsize=2,
    captionpos=t,
    frame=single, % <-- ADDED: Frame ensures background is applied
    framerule=0pt, % <-- ADDED: Hides the default frame
}

\lstdefinestyle{inlinecode}{
    backgroundcolor=\color{backcolour},   
    basicstyle=\ttfamily\footnotesize,
    keywordstyle=\color{blue},
    stringstyle=\color{codeorange},
    commentstyle=\color{codegreen},
}

% set styles
\lstset{style=mystyle,language=SQL}

% shortening för inline code
\lstMakeShortInline[style=inlinecode]`

% boilerplate code from https://www.guitex.org/home/images/ArsTeXnica/AT015/drawing-ER-diagrams-with-TikZ.pdf



% Title page
\title{\textbf{Homework 3}\\\Large{IS1500 group 74}}
\author{Isac Zobec \\ Linus Bredin \\ KTH Royal Institute of Technology\\  Computer Science (CDATE)}
\date{30/11, 2025}

% Header and footer
\pagestyle{fancy}
\fancyhead{}
\fancyhead[L]{\textbf{Database Technology}\\\textbf{DD1368}}
\fancyhead[R]{\textbf{Linus Bredin, Isac Zobec}\\ lbredin@kth.se izobec@kth.se}
\fancyfoot{}
\begin{document}

\maketitle
\thispagestyle{fancy}
\clearpage
\tableofcontents
\thispagestyle{fancy}

\clearpage
% Begin page numbers
\fancyfoot[C]{\thepage}
\pagenumbering{arabic}

\newpage

\section{Methodology}

\subsubsection*{Describe your methodology and how you performed the experiments.}

The performance of the program was measured by reading from the special control and status counter registers of the dtek board (`mcycle` through `mhpmcounter9`). The function `clearCsrRegs` was implemented to reset all the counters to 0, and the `printCsrRegs` function to print the values of all the counters to the console. In addition, to measure performance of a certain function/block over several calls, one variable to cumulatively store the value of each counter, and the functions `resetCsrCounterVars`, `incrementCsrCounterVars`, and `printCsrCounterVars`. were implemented. The first function resets all the variables to 0, the second adds the current values of the counters to the registers, and the third prints the variables' values to the terminal. The data was then copied to a spreadsheet to calculate the relevant metrics (IPC, cache hit/miss-rate).

For the rasterization/rendering part of the program, we cumulatively measured counters for one ''rendering-loop'', which is rendering all the pieces on the chessboard. We did this for just the part of the rendering function where the vertecies are transformed into clip space, and for the part where the triangles are rasterized onto the screen respectively.  

\section{Results}

Below are our results for the rendering part of the project.

\input{TriO1}

Note that the Miss rate is calculated as the amount of times an instruction fetch resulted in a D-cache miss (`mhpmcounter5`) divided by the amount of memory instructions retired (`mhpmcounter3`).

\subsubsection*{What is the IPC (Instruction-Per-Cycle) you achieve with your code on the DTEK-V board?}

For the triangle rasterization and vertex transformation parts, we get a low IPC of 0.35 and 0.38 respectively.

\subsubsection*{What is preventing you from reaching a peak (IPC=1) of your application? Motivate the question with empirical data from:}
\subsubsection*{Stalls due to delays in the instruction- and data-cache}

I-Cache misses are the main bottleneck in the program, as stalls due to I-Cache misses take up around 58.5\% of the cycles for both rendering parts. 

A likely reason for the I-Cache misses are that the rasterization code contains a tight rendering loop that is run very often, and that contains very many calls to different functions (math, color transformations, etc). 

\subsubsection*{Stalls due to delays in the arithmetic logic unit (ALU)}
\subsubsection*{Stalls due to data hazards}
\subsubsection*{What impact does compiling with no compiler optimizations (-O0) and heavy compiler optimizations (-O3) have on your application and on the above metric?}
\subsubsection*{A small discussion on what architectural improvements you would do to the processor and why you would do them in order for your basic project to run faster.}


% \section{P+}
% \centering\textit{\large{todo}}

% References
% \clearpage
% \addcontentsline{toc}{section}{References}
% \printbibliography

\end{document}
